\section{The determinant deficit at the original degree and the four-window scale}
\label{sec:original-degree-deficit}

Retain the complete original quartet and every constant of DMR1--37.
In particular the zero has full multiplicity $m\geq1$, and
\[
 \ell_k=1+k(m-1),\quad q=(k+1)^2\ell_k,\quad
 p=8m-\tfrac92,\quad M=21+4m,\quad
 \eta=(p-1)/p,\quad k\geq3.
 \tag{ODD1}
\]
The variable is still $S=k/2+iy$. Both coefficient forms retain
their original measures and masses. The fixed Gamma relative Gram
and its nonnegative determinant deficit are exactly
\[
 D_{k,N}=(H^{\sigma,k}_N)^{-1/2}H_{k,N}(H^{\sigma,k}_N)^{-1/2},
 \qquad \Xi_{k,N}=(N+1)\log A_k-\log\det D_{k,N}.
 \tag{ODD2}
\]
Every degree below is one of $q-1,q,2q-1,2q$. The present
calculation retains the full value of $q$ in the finite DMR bound.
It does not replace the source coordinate, polynomial, multiplicity
or observation.

For explicitness, the constants used from DMR30 are
\[
\begin{gathered}
 \alpha=\pi/2,\quad B=42+8m,\quad
 b=\frac{c_h}{C_\Gamma 2^{3B/2}},\quad
 C_r=1+\frac{\alpha+\log(M_\alpha/\vartheta_h)+B}{\log2},\\
 C_d=2C_r+M/3,\quad C_L=1+C_d,\quad
 C_E=6C_r^2+(2M+4)C_d/3,\\
 C_s=\left|\log\frac{C_h}{c_\Gamma b}+2\log M_\alpha\right|,
 \quad C_{\rm tail}=\frac{2C_h}{(p-1)M_1},\\
 T_k=\log(4m(C_L+2))+3\log(k+2),\\
 \vartheta_h=\int_{-1}^1w_h(y)\,dy,\quad
 M_1=\int_{-1}^1e^{\alpha y}w_h(y)\,dy,\quad
 M_\alpha=\int_{\mathbb R}e^{\alpha y}w_h(y)\,dy .
\end{gathered}\tag{ODD3}
\]
Here $C_h,c_h,c_\Gamma,C_\Gamma$ are the unchanged positive
envelope constants proved in TW and DMR2. In particular $M_1$
is the tilted finite mass and is not $\vartheta_h$. Every
quantity in ODD3 is finite, and the constants other than $T_k$
are independent of $k$ and of the truncation radius $R$.

Put
\[
 z_k=\frac{kT_k}{q}>0,\qquad
 R_k^*=\max(1,z_k^{-1/p}),\qquad
 B_k^*=(p+1)\log k+C_s+2MT_k.
 \tag{ODD4}
\]
These are explicit scalars on the unchanged source; $R_k^*$ is
an admissible radius for DMR29 and DMR32 at every $k\geq3$.
Define the following complete finite majorant:
\[
\begin{split}
 \Psi_k={}&\left(2+\frac1q\right)
 \left\{C_{\rm tail}\min(2^{1-p},z_k^\eta)
                       +\frac{B_k^*}{k}\right\}\\
 &+C_Ez_k+C_d\max(z_k,z_k^\eta)
                       +\frac{(g_0+2)C_d}{q},
 \qquad g_0=\log\sqrt{2\pi}.
\end{split}\tag{ODD5}
\]
Then the finite result at all four original degrees is
\[
 \boxed{\quad 0\leq\frac{\Xi_{k,N}}{kq}\leq\Psi_k.\quad}
 \tag{ODD6}
\]

Here is the proof with the degree factors retained. DMR29's
scalar and DMR32's exact error give, for every $R\geq1$,
\[
\begin{split}
 \frac{\Xi_{k,N}}{kq}\leq{}&
 \frac{N+1}{q}\left\{C_{\rm tail}(1+R)^{1-p}
                +\frac{(p+1)\log k+C_s}{k}\right\}\\
 &+\frac{2MN T_k}{kq}
       +\frac{kT_k}{q}(C_E+C_dR)+\frac{(g_0+2)C_d}{q}.
\end{split}\tag{ODD7}
\]
This follows by dividing the original finite inequality by $kq$;
no estimate for $k^2/(N+1)$ has been inserted. Since
$N/q\leq(N+1)/q\leq2+1/q$, the first and third logarithmic
terms combine to the $B_k^*/k$ term in ODD5. If $z_k\leq1$,
then $R_k^*=z_k^{-1/p}$ and
$z_kR_k^*=z_k^\eta$. If $z_k\geq1$, then $R_k^*=1$ and
$z_kR_k^*=z_k$. Thus in both cases
\[
 z_kR_k^*=\max(z_k,z_k^\eta),\qquad
 (1+R_k^*)^{1-p}\leq\min(2^{1-p},z_k^\eta).
 \tag{ODD8}
\]
For the second inequality, $R_k^*\geq1$ gives the first
bound. When $z_k\leq1$, the negative exponent gives
$(1+R_k^*)^{1-p}\leq(R_k^*)^{1-p}=z_k^\eta$.
When $z_k\geq1$, its value is $2^{1-p}\leq1\leq z_k^\eta$.
Substitution in ODD7 proves ODD6, including $z_k=1$ and $k=3$.

Because $T_k=O_h(\log k)$ and $q\geq(k+1)^2$, one has
$z_k\to0$. Formula ODD5 therefore proves
\[
 \frac{\Xi_{k,N}}{kq}
 =O_h\left(\frac{\log k}{k}
                   +\left(\frac{k\log k}{q}\right)^\eta\right).
 \tag{ODD9}
\]
Indeed $z_k\leq z_k^\eta$ for sufficiently large $k$ because
$0<\eta<1$, while $q^{-1}=O(k^{-2})$. This establishes
the displayed rate directly from a finite bound.

For $m\geq2$ the actual degree is
$q=(m-1)k^3+O_m(k^2)$ and $p\geq23/2$, so $\eta>1/2$.
Consequently
\[
 \frac{k}{\log k}\left(\frac{k\log k}{q}\right)^\eta
 =O_m\bigl(k^{1-2\eta}(\log k)^{\eta-1}\bigr)\longrightarrow0.
\]
Every term of ODD5 except $B_k^*/k$ vanishes after multiplication
by $k/\log k$: this follows from the last limit, from
$kz_k/\log k=O_m(k^{-1})$, and from
$k/(q\log k)\to0$. Since $T_k/\log k\to3$,
\[
 \boxed{\quad
 \limsup_{k\to\infty}\frac{\Xi_{k,2q}}{q\log k}
 \leq2\{p+1+6M\}=64m+245\qquad(m\geq2).
 \quad}\tag{ODD10}
\]
The equality retains $p+1=8m-7/2$ and $6M=126+24m$.
For $m=1$, the same proof instead retains
$q=(k+1)^2$ and $\eta=5/7$. ODD9 then gives
$\Xi_{k,N}/(kq)=O_h((\log k/k)^{5/7})$; ODD10 has
not been extended to that different degree regime.

The exact source and target maps in the comparison are as follows.
Let $J_N:\mathcal P_N\twoheadrightarrow E_k$ be the unchanged
remainder, and put $G_N=(J_NH_N^{-1}J_N^*)^{-1}$ for each
of the two original source Grams. The section is
$H_N^{-1}J_N^*G_N$; the kernel of $J_N$ is
$\chi\mathcal P_{N-q}$. For the original onto boundary
$\Lambda:E_k\to\mathcal B$ and its fixed kernel inclusion
$I_K:\mathcal K\hookrightarrow E_k$, retain
\[
 H_N^K=I_K^*G_NI_K,\quad
 Q_N^B=(\Lambda G_N^{-1}\Lambda^*)^{-1},\quad
 S_N^B=G_N^{-1}\Lambda^*Q_N^B .
 \tag{ODD11}
\]
The identities $\Lambda S_N^B=I$ and $I_K^*G_NS_N^B=0$
follow by multiplication. The complete positive determinant
decompositions QDC1--16 and QDC28--32, on these same maps,
give each original quotient deficit between zero and
$\Xi_{k,N}$ and split it into its two nonnegative kernel
and boundary deficits. The common ambient coefficient inclusion
$\mathcal P_N\hookrightarrow\mathcal P_{2q}$ gives
$\Xi_{k,N}\leq\Xi_{k,2q}$, as proved with its exact reference
isometry in DMR37. Taking the original signs $(+,+,-,-)$
on the four quotient logarithms gives $(-,-,+,+)$ on the
nonnegative deficits. Each positive or negative side has two
terms, hence
\[
 \max\{|\delta_k^\sigma|,|\Delta F_k^K|,|\Delta F_k^B|\}
 \leq2kq\Psi_k,
 \qquad
 \limsup_{k\to\infty}
 \frac{\max\{|\delta_k^\sigma|,|\Delta F_k^K|,|\Delta F_k^B|\}}
      {q\log k}\leq128m+490\quad(m\geq2).
 \tag{ODD12}
\]
The finite inequality holds for every $m\geq1$.

Finally retain HC5's original constant $D_h>0$, its full
nonnegative residual $\mathcal R_k$, and its exact four-volume
functional $\mathcal B_k$. Let $F_{\sigma,k}$ denote that same
functional for the fixed Gamma source. Their literal identities are
\[
 \mathcal R_k=F_{\sigma,k}+\delta_k^\sigma
                        -4q\log(D_hk)\geq0,
 \qquad
 \Delta_k^\Gamma=\delta_k^\sigma-\mathfrak T_k.
 \tag{ODD13}
\]
The first equality substitutes the exact definition of
$\delta_k^\sigma$ into HC5; the second is GCR27 with the
entire signed determinant $\mathfrak T_k$ of GCR26 retained.
Thus the proved finite original interval is
\[
 F_{\sigma,k}-4q\log(D_hk)-2kq\Psi_k
 \leq\mathcal R_k\leq
 F_{\sigma,k}-4q\log(D_hk)+2kq\Psi_k .
 \tag{ODD14}
\]
It also retains the independent exact constraint $\mathcal R_k\geq0$.
Since $4q\log(D_hk)/(q\log k)\to4$, ODD10--12 put the
multiple-zero arithmetic correction on this actual scale with
explicit constants. Formula ODD14 retains the unevaluated
fixed-Gamma functional $F_{\sigma,k}$. No sign for $\delta_k^\sigma$ or
$\mathfrak T_k$ has been inserted, and no RH endpoint follows
from the displayed enclosure alone.
